\documentclass[10pt,a4paper]{amsart}
\usepackage[margin=19mm]{geometry}
\usepackage{amsmath,amssymb,mathtools}
\usepackage[scr=rsfs,cal=euler]{mathalfa}
\usepackage{xfrac}
\usepackage[unicode,hidelinks,bookmarks=false]{hyperref}
\newcommand{\ie}{i.e.\ }
\newcommand{\cf}{cf.\ }
\newcommand{\resp}{respectively\ }
\newcommand{\Subsection}{\subsection}
\providecommand{\nobreakdash}{\nobreak\hbox{-}}
\setlength{\emergencystretch}{2em}
\multlinegap0pt
\usepackage{bm}
\usepackage{bbm}
\usepackage{dsfont}
\usepackage{xspace}
\usepackage{cancel}
\usepackage{mathtools}
\usepackage{amsthm}
\makeatletter
\@ifundefined{theorem}{\newtheorem{theorem}{Theorem}}{}
\@ifundefined{proposition}{\newtheorem{proposition}[theorem]{Proposition}}{}
\makeatother
\DeclareMathOperator*{\argmin}{arg\,min}
\title[Optimization of Bregman Variational Learning Dynamics]{Optimization of Bregman Variational Learning Dynamics}
\author{Jinho Cha, Youngchul Kim, Jungmin Shin, Jaeyoung Cho, Seon Jin Kim, Junyeol Ryu}
\date{Source: arXiv:2510.20227v1 (CC BY 4.0)}
\begin{document}
\maketitle

\noindent\emph{Source and licence.} Optimization of Bregman Variational Learning Dynamics. Version arXiv:2510.20227v1 (CC BY 4.0), distributed under the Creative Commons Attribution 4.0 International licence, as recorded in the arXiv licence metadata for this exact version and shown on the abstract page https://arxiv.org/abs/2510.20227v1. Full rights notice: licenses/arxiv-2510.20227-CC-BY-4.0.txt.\par\smallskip

\noindent\textit{Editorial scope.} Counted unit: the Proposition whose source label is \texttt{prop:bilevel} in the article (source0.tex lines 2488--2500, source label \texttt{prop:bilevel}), delivered together with the article's own complete original proof of it (source0.tex lines 2502--2564, one \texttt{proof} environment of 63 lines). No mathematics is written, repaired or re-derived: statement and proof are the article's own text line for line. Required context retained uncounted: eq:bilevel (lines 2477--2483). Every definition, display and label of those lines is retained unchanged. Omitted from this package and not redistributed: 1--2476, 2484--2487, 2501--2501, 2565--2749. Nothing inside a retained excerpt is omitted or altered: every retained body is delivered exactly as the article's lines are written, so no omission receipt of any kind accompanies this package. Citations: the retained excerpts carry no citation command at all, so no bibliography entry of the article is transcribed and no reference is redistributed. Reference adaptation: none was needed and none was made; every reference inside the delivered unit resolves inside it, so no unresolved reference or citation is produced. Printed slips kept literal: no printed word, symbol, hypothesis or number of the counted statement or of its proof was altered. Layout and class: the article's own class is \texttt{article} with packages including geometry, setspace, lmodern, fontenc, inputenc, amsmath, amssymb, bm, mathtools, booktabs, enumitem, graphicx, array, float; the wrapper is the lane's pinned \texttt{amsart} editorial wrapper at 10pt on A4, which restates the theorem-like environments the retained text needs and adds the article's own macro definitions that the retained text needs (\texttt{\textbackslash argmin}), with extra wrapper packages (bm, bbm, dsfont, xspace, cancel, mathtools). The delivered numbering is the unit's own. Assets: no retained span contains a figure environment, a graphics-inclusion command, a PDF-page inclusion, a table or a TikZ picture; no image file is copied into this package, the package declares no asset dependency and no image file is redistributed with it. Licence: version arXiv:2510.20227v1 of this article is distributed under the Creative Commons Attribution 4.0 International licence, as recorded in the arXiv licence metadata for this exact version and shown on the abstract page https://arxiv.org/abs/2510.20227v1. A full rights notice is delivered at licenses/arxiv-2510.20227-CC-BY-4.0.txt. Characters: every retained excerpt is pure ASCII, so the delivered engine, XeLaTeX with its default Latin Modern text font, reports no missing character.
\begin{equation}
\min_{x}\; C(x,q)+\lambda\,\mathrm{Risk}(x,q)
\quad
\text{s.t.}\quad
q\in\argmin_{r\in\Theta}\{f_t(r)+D_\psi(r\|p)\}.
\label{eq:bilevel}
\end{equation}

\begin{proposition}[KKT-Based Single-Level Reformulation]
\label{prop:bilevel}
Embedding the optimality condition of the lower-level BVLD problem 
yields the equivalent single-level convex program:
\begin{equation}
\min_{x,q}\;
C(x,q)+\lambda\,\mathrm{Risk}(x,q)
\quad
\text{s.t.}\quad
0\in\nabla f_t(q)+\nabla\psi(q)-\nabla\psi(p)+N_\Theta(q).
\label{eq:singlelevel}
\end{equation}
\end{proposition}

\begin{proof}
We proceed in four steps.

\emph{(i) Well-posedness and uniqueness of the lower-level BVLD subproblem.}
Fix $p\in\Theta$. Consider the lower problem 
\[
\min_{r\in\Theta}\; f_t(r)+D_\psi(r\|p).
\]
By convexity of $f_t$ and strong convexity of $D_\psi(\cdot\|p)$ in $r$, 
the objective is strongly convex on $\Theta$.  
If $\Theta$ is unbounded, the strong convexity of $D_\psi$ ensures coercivity; 
hence a unique minimizer $q^\star\in\Theta$ exists.

\emph{(ii) Necessary and sufficient optimality (KKT/monotone inclusion).}
Since the objective is convex and $C^1$, and the feasible set $\Theta$ is closed and convex,
first-order optimality is necessary and sufficient:
\[
0\in \nabla f_t(q^\star)+\nabla_q D_\psi(q^\star\|p)+N_\Theta(q^\star).
\]
Using the Bregman identity $\nabla_q D_\psi(q\|p)=\nabla\psi(q)-\nabla\psi(p)$, 
we obtain
\begin{equation}\label{eq:bvld-kkt}
0\in\nabla f_t(q^\star)+\nabla\psi(q^\star)-\nabla\psi(p)+N_\Theta(q^\star).
\end{equation}
Conversely, any $q$ satisfying \eqref{eq:bvld-kkt} 
is the unique minimizer of the lower problem, 
so \eqref{eq:bvld-kkt} characterizes the feasible set of the lower-level argmin exactly.

\emph{(iii) Equivalence of the bilevel and single-level feasible pairs.}
The feasible set of the bilevel model
\[
\big\{(x,q):\; q\in\argmin_{r\in\Theta}\{f_t(r)+D_\psi(r\|p)\}\big\}
\]
coincides with
\[
\big\{(x,q):\; 0\in\nabla f_t(q)+\nabla\psi(q)-\nabla\psi(p)+N_\Theta(q)\big\},
\]
by (ii).  
Therefore, the constraint in~\eqref{eq:bilevel} 
is equivalently expressed by the KKT/normal-cone inclusion 
in~\eqref{eq:singlelevel}.  
No spurious solutions are introduced and none are lost.

\emph{(iv) Convexity and optimality preservation under embedding.}
Under the standing assumptions, 
the objective $C(x,q)+\lambda\,\mathrm{Risk}(x,q)$ is convex in $(x,q)$.  
The set-valued map
\[
q\mapsto\nabla f_t(q)+\nabla\psi(q)-\nabla\psi(p)+N_\Theta(q)
\]
is \emph{maximally monotone}: $\nabla f_t$ and $\nabla\psi$ are monotone 
(with $\nabla\psi$ $\mu$-strongly monotone), and $N_\Theta$ is maximally monotone.  
Hence the feasible set
\(
\{q:\,0\in\nabla f_t(q)+\nabla\psi(q)-\nabla\psi(p)+N_\Theta(q)\}
\)
is closed and convex.  
Consequently, the single-level formulation~\eqref{eq:singlelevel} 
is a convex optimization problem.  
Since both formulations share the same objective and feasible pairs, 
their optimal solutions coincide.  
% \qed
\end{proof}

\end{document}
